\section{Discussion and lessons learned}
\label{sec:discussion}

The evaluation distinguishes document discovery from evidence-supported
answering. The retrieval-only matrix shows that document-level recall alone
does not account for answer-layer behavior: the same retrieved document can
support a citation, an explicit partial answer, or an honest refusal,
depending on the evidence available to the agent. A knowledge-base system
should therefore expose both retrieval outcomes and the evidence path that
produced an answer. The three contracts of Section~\ref{sec:system} name that
path, and the lessons below follow it from ingest to answer.

Verdict placement is the first lesson. On identical candidate pools the same
judge scores hold fusion-quality retrieval when they act as an
admission filter (Hit@5 0.94) and lose the gold when they order the list
(0.44), because same-domain near-misses inflate the scores until the ordering
carries almost no gold information; a second, generic verifier fails the same
way, so replacing the checkpoint is not the repair. What the gate buys, on the
measured corpora, is fail-closed provenance rather than ranking power: every
returned document was fetched and judged, and the verdict travels with the
evidence. That provenance is inspectable precisely because the verifier sits
server-side, outside any single generation loop.

Abstention is the second lesson. A similarity score retrieves topically
related material even when no document answers the query, so refusal has to
be enforced in the reading and answer layers, where the system can state
what it searched and where its coverage ends. The adversarial battery
recorded no fabricated citations across eighteen final arms, but its
single-run design and one body-level partial answer limit that result, and
a binary found/refused gate cannot represent an answer that cites applied
evidence while disclosing that a dedicated tutorial is absent.

Tool transport is the third lesson: the oversized-catalog incident
reproduced a false refusal in which the document was present but the agent
received only a truncated preview, and server-side budgets, grouped catalog
views and pagination repaired the tested path on rerun --- evidence for
that boundary, not a guarantee for every model host or larger catalog.

The five questions of Section~\ref{sec:rq} therefore close as follows. RQ1:
the contracts bind where their checks are server-side mechanisms; the stale
index of Section~\ref{sec:e2} was caught by the three-source audit, while the
truncated catalog of Section~\ref{sec:sysarms} showed that an enforcement
layer has transport limits of its own, closed by budget-guarded pagination.
RQ2: the acquisition gates and content routing behaved as specified on
deployed collections, though 21 of 100 submitted papers
never became retrievable and routing is demonstrated as coverage, not
accuracy. RQ3: at this scale the lanes separate on checkability, citation
discipline and cost, not on document recall. RQ4: content verification earns
its keep as an admission filter with fail-closed provenance, fails as a
re-ranker under same-domain confuser density, and survives no checkpoint
substitution; its calibration remains untested. RQ5: abstention is enforced
at the answer layer, latency is turn-dominated, and the delivery surface
moves cost by an order of magnitude, with tool-output budgets a constraint
on the enforcement layer itself.

More broadly, parsers, registries, indices and answer agents should expose
post-conditions checkable independently of a success envelope; catalogs and
graphs aid navigation but do not substitute for content-level evidence,
consistent with provenance and auditability principles in data systems
\cite{missier-prov-2013,cheney-provenance-2009} and with the fallibility of
learned judges \cite{zheng-llmjudge,panickssery-selfpref}. The evidence
remains bounded by the tested corpora, developer-authored questions, one
workstation and few end-to-end runs; larger independently curated
collections, repeated agent trials, operator studies of not-found reports
and crash-recovery tests are needed before generalizing these contracts to
broader deployment settings.
